\documentclass[12pt, a4paper]{article}

% ── Packages ────────────────────────────────────────────────────────────────
\usepackage[T1]{fontenc}
\usepackage[utf8]{inputenc}
\usepackage{lmodern}
\usepackage{microtype}
\usepackage{amsmath, amssymb, amsthm}
\usepackage{mathtools}
\usepackage{enumitem}
\usepackage{geometry}
\usepackage{hyperref}
\usepackage{booktabs}
\usepackage{array}
\usepackage{xcolor}
\usepackage{listings}
\usepackage{fancyhdr}
\usepackage{titlesec}
\usepackage{tcolorbox}
\tcbuselibrary{theorems, skins, breakable}

\geometry{margin=2.5cm}

% ── Theorem environments ─────────────────────────────────────────────────────
\newtheorem{theorem}{Theorem}[section]
\newtheorem{lemma}[theorem]{Lemma}
\newtheorem{corollary}[theorem]{Corollary}
\newtheorem{proposition}[theorem]{Proposition}
\theoremstyle{definition}
\newtheorem{definition}[theorem]{Definition}
\newtheorem{remark}[theorem]{Remark}

% ── Code listings ─────────────────────────────────────────────────────────
\definecolor{codegray}{rgb}{0.95,0.95,0.95}
\definecolor{codeblue}{rgb}{0.1,0.2,0.6}
\definecolor{darkgreen}{rgb}{0.0,0.45,0.0}
\lstset{
  backgroundcolor=\color{codegray},
  basicstyle=\ttfamily\small,
  keywordstyle=\color{codeblue}\bfseries,
  breaklines=true,
  frame=single,
  framerule=0pt,
  rulecolor=\color{gray!30},
  xleftmargin=6pt,
  xrightmargin=6pt,
}

% ── Lean listings ─────────────────────────────────────────────────────────
\lstdefinelanguage{Lean4}{
  keywords={theorem,lemma,def,axiom,by,have,obtain,exact,intro,rw,push_cast,
            ring,decide,linarith,nlinarith,rcases,calc,simp,fin_cases,revert,
            import,infix,infixl,infixr,where,fun,if,then,else,let,in,
            return,do,match,with,from,show,suffices,apply,refine,constructor,
            left,right,exfalso,contradiction,norm_num,norm_cast,exact_mod_cast,
            native_decide,omega},
  sensitive=true,
  morecomment=[l]{--},
  morecomment=[s]{/-}{-/},
  morestring=[b]",
  keywordstyle=\color{codeblue}\bfseries,
  commentstyle=\color{gray}\itshape,
}

% ── Macros ──────────────────────────────────────────────────────────────────
\newcommand{\Z}{\mathbb{Z}}
\newcommand{\Q}{\mathbb{Q}}
\newcommand{\N}{\mathbb{N}}
\newcommand{\F}{\mathbb{F}}

% ── Header / Footer ─────────────────────────────────────────────────────────
\pagestyle{fancy}
\fancyhf{}
\rhead{\small\thepage}
\lhead{\small No Integer Solutions to $y^2z + yz^2 = x^3 + x^2 + 3x - 1$}
\renewcommand{\headrulewidth}{0.4pt}

% ── tcolorbox styles ────────────────────────────────────────────────────────
\tcbset{
  mythm/.style={
    enhanced, colback=blue!4!white, colframe=blue!40!black,
    fonttitle=\bfseries, breakable,
  },
  mylemma/.style={
    enhanced, colback=green!4!white, colframe=darkgreen,
    fonttitle=\bfseries, breakable,
  },
  myobs/.style={
    enhanced, colback=orange!6!white, colframe=orange!60!black,
    fonttitle=\bfseries, breakable,
  },
}

% ────────────────────────────────────────────────────────────────────────────
\begin{document}

\title{\bfseries No Integer Solutions to\\[4pt]
  $y^2 z + y z^2 = x^3 + x^2 + 3x - 1$}
\author{}
\date{April 2026}
\maketitle

\begin{abstract}
We study the Diophantine equation
\[
  y^2 z + y z^2 \;=\; x^3 + x^2 + 3x - 1, \qquad (x,y,z)\in\Z^3.
\]
Equivalently, writing $yz(y+z)$ for the left side, we ask for which
integer triples the product of $y$, $z$, and their sum equals the cubic
$x^3+x^2+3x-1$.

We prove:
\begin{enumerate}
  \item $x$ must be odd (parity argument, \S\ref{sec:parity}).
  \item For every odd $x$, the right side is exactly divisible by $4$
        but not $8$ (2-adic analysis, \S\ref{sec:twadic}).
  \item For $x \equiv 1,\,5,\,9,\,13\pmod{14}$ the right side is
        $\equiv 3$ or $4\pmod 7$, a value never achieved by $yz(y+z)$
        (mod-7 obstruction, \S\ref{sec:mod7}).
\end{enumerate}
This yields a rigorous elementary proof that no solution exists for
the four residue classes $x\equiv 1,5,9,13\pmod{14}$, which together
account for $4/7$ of all odd integers.

For the remaining classes $x\equiv 3,\,7,\,11\pmod{14}$, the equation
is locally solvable at every prime (no further modular obstruction
exists), and an exhaustive divisor-enumeration search confirms no
solutions for $|x|\le 10{,}000$.  Completing the proof for these
classes requires an elliptic-curve descent argument.

\medskip\noindent\textbf{Conjecture.}  The equation has no integer solutions.
\end{abstract}

\tableofcontents

% ============================================================
\section{Setup and notation}
% ============================================================

The equation
\begin{equation}\label{eq:main}
  y^2 z + y z^2 = x^3 + x^2 + 3x - 1
\end{equation}
factors on the left as $yz(y+z)$.  We write
\[
  f(x) := x^3 + x^2 + 3x - 1
\]
and
\[
  g(y,z) := yz(y+z).
\]
A triple $(x,y,z)\in\Z^3$ is a \emph{solution} if $g(y,z)=f(x)$.

In the substitution $s = y+z$, $p = yz$, the condition becomes $sp = f(x)$,
with $y,z$ real iff $s^2-4p\ge0$ and integer iff additionally $s^2-4p$ is
a perfect square and $s^2-4p\equiv s^2\pmod{4}$.

% ============================================================
\section{Parity: \texorpdfstring{$x$}{x} must be odd}
\label{sec:parity}
% ============================================================

\begin{lemma}\label{lem:parity}
  For any $(x,y,z)\in\Z^3$ satisfying \eqref{eq:main}, $x$ must be odd.
\end{lemma}

\begin{proof}
We show $g(y,z)$ is always even. Fix $y,z\in\Z$ and consider three cases:
\begin{itemize}
  \item \emph{Both $y,z$ even.} Then $yz\equiv 0\pmod{2}$, so $g(y,z)\equiv0$.
  \item \emph{Exactly one even.} Then $yz\equiv0\pmod{2}$, so $g(y,z)\equiv0$.
  \item \emph{Both $y,z$ odd.} Then $y+z\equiv0\pmod{2}$, so again
        $g(y,z)=yz(y+z)\equiv0\pmod{2}$.
\end{itemize}
In all cases $g(y,z)\equiv 0\pmod{2}$.

On the right, $f(x)\equiv x^3+x^2+3x-1\equiv x+x+x+1=3x+1\equiv x+1\pmod{2}$.
So $f(x)$ is even iff $x$ is odd. The claim follows.
\end{proof}

% ============================================================
\section{2-adic structure}
\label{sec:twadic}
% ============================================================

\begin{lemma}\label{lem:twadic}
  For every odd integer $x$, $\nu_2(f(x))=2$, i.e.\ $4\mid f(x)$ but
  $8\nmid f(x)$.
\end{lemma}

\begin{proof}
Write $x = 2k+1$.  Then
\begin{align*}
  f(2k+1) &= (2k+1)^3 + (2k+1)^2 + 3(2k+1) - 1 \\
  &= (8k^3 + 12k^2 + 6k + 1) + (4k^2 + 4k + 1)
     + (6k + 3) - 1 \\
  &= 8k^3 + 16k^2 + 16k + 4 \\
  &= 4\underbrace{(2k^3 + 4k^2 + 4k + 1)}_{=: q(k)}.
\end{align*}
We claim $q(k)$ is odd for every integer $k$.  Indeed,
$q(k)\equiv 2k^3+4k^2+4k+1\equiv 1\pmod{2}$ (since the terms $2k^3$,
$4k^2$, $4k$ are all even).  Hence $\nu_2(f(x))=\nu_2(4q(k))=2$.
\end{proof}

\begin{remark}
The lemma implies that in any solution we must have $\nu_2(g(y,z))=2$,
which constrains $(y,z)$ modulo $8$ to the two cases in
Lemma~\ref{lem:twadic-yz} below.
\end{remark}

\begin{lemma}\label{lem:twadic-yz}
$\nu_2(yz(y+z))=2$ if and only if, up to swapping $y$ and $z$:
\begin{enumerate}[(a)]
  \item $y$ and $z$ are both odd, and $y+z\equiv 4\pmod{8}$, or
  \item $y\equiv4\pmod{8}$ and $z$ is odd.
\end{enumerate}
\end{lemma}

\begin{proof}
Direct case analysis on $\nu_2(y)$, $\nu_2(z)$, and $\nu_2(y+z)$,
using $\nu_2(ab)=\nu_2(a)+\nu_2(b)$. The key observations are:
if both $y,z$ are even then $4\mid yz$, so $\nu_2(g)\ge 4$;
if one is $\equiv 2\pmod4$ and the other is odd then $\nu_2(yz)=1$
and $\nu_2(y+z)=0$, giving $\nu_2(g)=1$. Case~(a) and (b) cover
the remaining situations yielding $\nu_2(g)=2$.
\end{proof}

% ============================================================
\section{The mod-7 obstruction}
\label{sec:mod7}
% ============================================================

\begin{tcolorbox}[myobs, title={Key Theorem}]
\textbf{Theorem.}
$g(y,z) = yz(y+z)$ satisfies $g(y,z)\not\equiv 3\pmod7$ and
$g(y,z)\not\equiv 4\pmod7$ for every $(y,z)\in\Z^2$.
\end{tcolorbox}

\begin{proof}
It suffices to check all $(y,z)\in(\Z/7\Z)^2$.  The non-zero values
are obtained for $y,z\not\equiv0\pmod7$ and $y+z\not\equiv0\pmod7$.
The complete table of $yz(y+z)\pmod7$ for such pairs is:

\begin{center}
\begin{tabular}{c|cccccc}
\toprule
$y\backslash z$ & 1 & 2 & 3 & 4 & 5 & 6 \\\midrule
1 & 2 & 6 & 5 & 6 & 2 & 0 \\
2 & 6 & 1 & 5 & 6 & 0 & 5 \\
3 & 5 & 5 & 5 & 0 & 1 & 1 \\
4 & 6 & 6 & 0 & 2 & 5 & 2 \\
5 & 2 & 0 & 1 & 5 & 5 & 6 \\
6 & 0 & 5 & 1 & 2 & 6 & 1 \\
\bottomrule
\end{tabular}
\end{center}

All entries lie in $\{0,1,2,5,6\}$; neither $3$ nor $4$ appears.
\end{proof}

\begin{corollary}\label{cor:mod7}
For the following residue classes modulo $14$, the equation \eqref{eq:main}
has no integer solutions.
\begin{center}
\begin{tabular}{ccc}
\toprule
$x \pmod{14}$ & $f(x)\pmod7$ & Obstruction \\\midrule
$1$  & $4$ & $4\notin\{0,1,2,5,6\}$ \\
$5$  & $3$ & $3\notin\{0,1,2,5,6\}$ \\
$9$  & $3$ & $3\notin\{0,1,2,5,6\}$ \\
$13$ & $3$ & $3\notin\{0,1,2,5,6\}$ \\\bottomrule
\end{tabular}
\end{center}
Combined with Lemma~\ref{lem:parity} (only odd $x$ can occur), this
eliminates \emph{four of the seven odd residue classes modulo~$14$},
i.e.\ $4/7$ of all odd integers.
\end{corollary}

\begin{proof}
Compute $f(x)\pmod7$ for each $x\in\{1,5,9,13\}$ (representative values
for the stated congruence classes):
\begin{align*}
  f(1)  &= 4,\quad &f(5) &= 164 = 23\cdot7+3\equiv3,\\
  f(9)  &= 836 = 119\cdot7+3\equiv3,\quad &f(13) &= 2404 = 343\cdot7+3\equiv3.
\end{align*}
Since $g(y,z)\ne3$ and $g(y,z)\ne4$ over $\Z/7\Z$, the equation is
insoluble modulo~7 for these $x$.
\end{proof}

% ============================================================
\section{Local solvability of the surviving cases}
\label{sec:local}
% ============================================================

The surviving odd residue classes are $x\equiv 3,7,11\pmod{14}$.

\begin{proposition}
For $x\equiv3,\,7,\,11\pmod{14}$ and every prime $p$, the equation
$g(y,z)\equiv f(x)\pmod{p}$ is solvable.
\end{proposition}

\begin{proof}[Sketch]
For a prime $p\ge5$, the curve $C_N: yz(y+z)=N$ over $\F_p$ is a smooth
genus-1 cubic whenever $N\not\equiv0\pmod{p}$ (the only singular point
at $y=z=0$ satisfies $N\equiv0$).   By Hasse's theorem,
$\lvert\#C_N(\F_p)-(p+1)\rvert\le2\sqrt{p}$, so $\#C_N(\F_p)\ge(\sqrt p-1)^2>0$
for $p\ge4$.  When $N\equiv0\pmod p$ one checks directly using $y=0$ or $z=0$.
For $p=2,3$ an explicit check confirms solvability.

A computer search (`modular\_analysis.py`) verifies local solvability for
all primes $p\le113$.
\end{proof}

\begin{remark}
A systematic search over all single moduli $m\le1000$ and all
pairs $(7,m)$ with $m\le1000$ finds no modular combination that
eliminates the surviving classes. The surviving cases are
\emph{not} obstructed by any finite set of congruences, so a
complete proof must use a global argument.
\end{remark}

% ============================================================
\section{Computational verification}
\label{sec:computation}
% ============================================================

\begin{proposition}\label{prop:computational}
The equation \eqref{eq:main} has no integer solutions with $|x|\le10{,}000$.
\end{proposition}

\begin{proof}
The script \texttt{brute\_force\_search.py} performs an exhaustive search.
For each odd $x\in[-10{,}000,\,10{,}000]$:
\begin{enumerate}
  \item Compute $N=f(x)$.
  \item Find all integer divisors $s$ of $N$.
  \item For each $s$, set $p=N/s$ and compute $\Delta=s^2-4p$.
  \item Check whether $\Delta\ge0$ is a perfect square $k^2$ with
        $s+k\equiv0\pmod{2}$.
  \item If so, $(y,z)=\bigl((s+k)/2,\,(s-k)/2\bigr)$ is an integer
        solution; record it.
\end{enumerate}
Every integer solution $(y,z)$ to $yz(y+z)=N$ gives a divisor
$s=y+z$ of $N$, so the algorithm is complete: it misses no solution.

The result is zero solutions.
\end{proof}

% ============================================================
\section{Elliptic curve structure of the surviving fibers}
\label{sec:ec}
% ============================================================

For fixed $x$ with $N:=f(x)\ne0$, the equation $g(y,z)=N$ defines a
plane cubic in $(y,z)$.  Completing the square via $z=t-y$ gives
$y(t-y)\cdot t=N$, i.e.\
\[
  t y^2 - t^2 y + N = 0,
\]
a conic parametrised by $t$.  Equivalently, with $s=y+z$, $p=yz$:
\[
  E_N\colon sp = N, \quad s^2-4p = k^2 \quad (k\in\Z).
\]
After the substitution $s=u$, $p=N/u$, one obtains
\[
  u^2 - 4N/u = k^2 \quad\Longleftrightarrow\quad u^3 - u k^2 = 4N.
\]
Setting $w=k$ this is the Weierstrass model of an elliptic curve.
By Siegel's theorem, each $E_N$ has only finitely many integer points.
However, the integer-point sets $E_{f(x)}(\Z)$ must be shown to be
empty simultaneously for all surviving $x$.

A natural approach for the remaining cases $x\equiv3,7,11\pmod{14}$
is a 2-descent on the family $\{E_{f(x)}\}$, exploiting the uniform
congruence structure within each residue class.

% ============================================================
\section{Lean~4 formalisation}
\label{sec:lean}
% ============================================================

The file \texttt{NonExistenceProof.lean} contains a Lean~4 / Mathlib~4
formalisation built against Mathlib v4.21.0.

\textbf{Sorry count: 1} — the single remaining \texttt{sorry} covers
the non-existence of solutions for $x\equiv3,\,7,\,11\pmod{14}$.

\medskip
The following are \textbf{fully proved without any \texttt{sorry}}:

\begin{center}
\begin{tabular}{lll}
\toprule
Lean name & Statement & Method\\\midrule
\texttt{lhs\_always\_even} & $g(y,z)\equiv0\pmod2$ for all $y,z$ & \texttt{decide} on \texttt{ZMod 2}\\
\texttt{rhs\_even\_iff\_x\_odd} & $f(x)\equiv0\pmod2\iff x\text{ odd}$ & \texttt{decide}\\
\texttt{x\_odd} & Any solution has $x$ odd & parity lemmas\\
\texttt{nu2\_rhs\_eq\_two} & $\nu_2(f(x))=2$ for odd $x$ & \texttt{omega} + algebra\\
\texttt{g\_mod7\_image} & $yz(y+z)\pmod7\in\{0,1,2,5,6\}$ & \texttt{decide} on \texttt{ZMod 7}\\
\texttt{f\_mod7\_bad\_classes} & $f(x)\pmod7\in\{3,4\}$ for $x\equiv1,5,9,13\pmod{14}$ & \texttt{decide}\\
\texttt{no\_solution\_bad\_classes} & No solution for $x\equiv1,5,9,13\pmod{14}$ & above two lemmas\\
\bottomrule
\end{tabular}
\end{center}

\begin{lstlisting}[language=Lean4,caption={Key excerpt from NonExistenceProof.lean}]
-- yz(y+z) mod 7 never equals 3 or 4
lemma g_mod7_image (y z : ZMod 7) :
    y * z * (y + z) ≠ 3 ∧ y * z * (y + z) ≠ 4 := by decide

-- f(x) ≡ 3 or 4 (mod 7) for x ≡ 1,5,9,13 (mod 14)
lemma f_mod7_bad_classes (x : ZMod 14) 
    (hx : x = 1 ∨ x = 5 ∨ x = 9 ∨ x = 13) :
    (x.val^3 + x.val^2 + 3*x.val - 1 : ZMod 7) = 3 ∨
    (x.val^3 + x.val^2 + 3*x.val - 1 : ZMod 7) = 4 := by
  rcases hx with rfl | rfl | rfl | rfl <;> decide
\end{lstlisting}

% ============================================================
\section{Summary}
% ============================================================

\begin{tcolorbox}[mythm, title={Main Result (partial proof)}]
\textbf{Theorem.}
The equation $y^2z + yz^2 = x^3 + x^2 + 3x - 1$ has no integer solutions
$(x,y,z)$ with $x \equiv 1, 5, 9,$ or $13\pmod{14}$.
\medskip

\noindent\textbf{Proof summary.}
By Lemma~\ref{lem:parity}, $x$ must be odd.
By Theorem~\ref{sec:mod7} and Corollary~\ref{cor:mod7},
$yz(y+z)\not\equiv 3, 4\pmod7$, yet $f(x)\equiv3$ or $4\pmod7$ for
$x\equiv1,5,9,13\pmod{14}$.  This gives the contradiction.
\end{tcolorbox}

\begin{tcolorbox}[myobs, title={Computational Evidence}]
An exhaustive divisor-enumeration search over $|x|\le 10{,}000$ finds
\textbf{zero solutions} to the full equation.  Combined with the
elementary proof above, this gives very strong evidence that the
equation has \emph{no} integer solutions whatsoever.  A complete proof
for $x\equiv3,7,11\pmod{14}$ remains open pending a suitable
descent argument.
\end{tcolorbox}

\end{document}
